\section{Evaluation protocol and controls}
\label{sec:protocol}

The protocol is where most of our disagreements with earlier numbers originate, so we state it
in full.

\subsection{Splits}

\paragraph{Respondent-disjoint (primary).} Every event of a respondent, and of a household where
the identifier encodes one, is assigned to a single partition. No respondent-specific effect is
fitted for a held-out respondent; the fallback is the zero effect. This is the deployment
question of predicting for a person never seen before.

\paragraph{Time-respecting (secondary).} We first verify that the ordering field is calendar
time. Of our three travel datasets only one passes: \dataset{LPMC} carries $3{,}371$ distinct
timestamps over three years, \dataset{Swissmetro}'s ordering field is the presentation index of a
stated-preference questionnaire, and \dataset{Optima}'s dates were manufactured by a preparation
script and take four distinct values. We therefore report the time-respecting protocol on
\dataset{LPMC} only, with forward-only blocks, and record the other two as unsupported with that
reason rather than inventing timestamps.

\paragraph{The split the earlier results used.} A chronological split \emph{within} respondents
places the same people in training and test. We keep it as a named reproduction configuration,
never as a primary result, and \Cref{tab:protocol} shows why it matters.

\subsection{Nested calibration}

The two scalars in \Cref{eq:mixture} are fitted by L-BFGS on predictions that neither channel
trained on. We use five outer folds over the development partition; for each fold the entire
pipeline -- both channels -- is rebuilt on the complement and predicts the held-out fold. The
scalars are then fitted on the pooled out-of-fold predictions and the stages are refitted on the
full development partition before the locked test partition is scored once. Folds hold out whole
respondents when the test regime is respondent-disjoint and hold out events when it shares
respondents, so the fold matches the regime the weight will be used in. The temperature is
applied exactly once, which we enforce with a test, and the numeric-only baseline receives its
own temperature fitted on the same held-out data with the same budget.

Fitting the weight on a single validation split, as the original pipeline did, gives a weight
that differs from the test-optimal value by up to $0.3$; out-of-fold fitting removes the largest
of the resulting losses (\Cref{app:calibration}).

\subsection{A hierarchy of controls}
\label{sec:controls}

Each control removes one explanation. In increasing strength:

\begin{description}
\item[Second numeric channel.] A parameter-budget-matched model with no text, mixed through the
      same calibration. Isolates ensembling.
\item[Shuffled consequences.] Sentences from a different event, matched on dataset, alternative
      and axis, with donors drawn only from the fitting partition and never from the recipient
      respondent. Isolates event-specific content from generic content.
\item[Identity-only text.] A sentence naming the alternative and nothing else. Isolates
      alternative identity.
\item[Random embeddings.] Seeded unit vectors with the encoder's normalisation. Isolates the
      embedding geometry.
\item[Templated facts.] Deterministic sentences built from exactly the fields the generator is
      allowed to see, for example \emph{``Recorded out-of-pocket cost for car: 1.18 CHF.''}
      Isolates what the language adds beyond restating the recorded quantities.
\item[Uniform floor.] The language channel replaced by the uniform distribution at the same
      fitted $\pi$.
\end{description}

The last is the one we did not find in the literature and the one that changes a conclusion.
Because $p_i(y) \geq (1-\pi)\, q_i(y)$ for the mixture in \Cref{eq:mixture},
\begin{equation}
-\log p_i(y) \;\leq\; -\log q_i(y) \;-\; \log(1-\pi),
\label{eq:floor}
\end{equation}
so mixing in \emph{any} distribution bounds the loss on events where the conventional model is
confidently wrong. At $\pi = 0.3$ that bound is $0.36$ nats per event, and it is available to a
channel that predicts at chance. The first four controls are all trained models producing
non-uniform predictions and so do not isolate this; the uniform floor does.

\subsection{Metrics and intervals}

The primary metric is mean per-event negative log-likelihood; Brier score, top-1 accuracy and a
fifteen-bin top-label calibration error are also recorded. Differences are paired per event as
$\mathrm{NLL}_{\text{baseline}} - \mathrm{NLL}_{\text{variant}}$, so positive favours the variant.
Intervals are $2{,}000$-draw bootstraps resampling \emph{respondents} (households on
\dataset{LPMC}) with all of a cluster's events moving together and the same resampled clusters
across variants; they condition on the trained models and the observed split. Five prespecified
contrasts form a Holm-corrected confirmatory family; the rest is exploratory, and an interval
containing zero is never read as equivalence.
